\subsection{\texorpdfstring{正合列与 Ext 元素之间的关系 \(_{A}\) (M, N)}{正合列与 Ext 元素之间的关系 \_\{A\} (M, N)}}

定理 1. --- 设 n 为整数 ≥1，M 和 N 为两个 A-模。

\begin{enumerate}
\def\labelenumi{\alph{enumi})}
\item
  Ext \(_{A}^{n}\) (M, N) 的每个元素都是某个正合列的类（X，第 117 页，定义 1）。
\item
  设 \(0 \to \mathbf{N} \xrightarrow{f_{n+1}} \mathbf{R}_n \xrightarrow{f_n} \ldots \to \mathbf{R}_1 \xrightarrow{f_1} \mathbf{M} \to 0\)
\end{enumerate}

和

\[
0 \rightarrow \mathrm{N} \xrightarrow {f _ {n + 1} ^ {\prime}} \mathrm{R} _ {n} ^ {\prime} \xrightarrow {f _ {n} ^ {\prime}} \dots \rightarrow \mathrm{R} _ {1} ^ {\prime} \xrightarrow {f _ {1} ^ {\prime}} \mathrm{M} \rightarrow 0
\]

为两个正合列， \(\theta\) 与 \(\theta'\) 为其相伴的类。以下条件等价：

\begin{enumerate}
\def\labelenumi{(\roman{enumi})}
\item
  \(\theta = \theta^{\prime}\) ;
\item
  存在一个具有正合行的交换图：
\end{enumerate}

\[
\begin{array}{c} 0 \longrightarrow \mathrm{N} \xrightarrow {f _ {n + 1}} \mathrm{R} _ {n} \longrightarrow \dots \longrightarrow \mathrm{R} _ {1} \xrightarrow {f _ {1}} \mathrm{M} \longrightarrow 0 \\ 0 \longrightarrow \mathrm{N} \longrightarrow \mathrm{R} _ {n} ^ {\prime \prime} \longrightarrow \dots \longrightarrow \mathrm{R} _ {1} ^ {\prime} \longrightarrow \mathrm{M} \longrightarrow 0 \\ 0 \longrightarrow \mathrm{N} \xrightarrow {f _ {n + 1} ^ {\prime}} \mathrm{R} _ {n} ^ {\prime} \longrightarrow \dots \longrightarrow \mathrm{R} _ {1} ^ {\prime} \xrightarrow {f _ {1} ^ {\prime}} \mathrm{M} \longrightarrow 0; \end{array}
\]

\begin{enumerate}
\def\labelenumi{(\roman{enumi})}
\setcounter{enumi}{2}
\tightlist
\item
  存在一个具有正合行的交换图：
\end{enumerate}

\[
\begin{array}{c} 0 \longrightarrow \mathrm{N} \xrightarrow {f _ {n + 1}} \mathrm{R} _ {n} \longrightarrow \dots \longrightarrow \mathrm{R} _ {1} \xrightarrow {f _ {1}} \mathrm{M} \longrightarrow 0 \\ 0 \longrightarrow \mathrm{N} \longrightarrow \mathrm{R} _ {n} ^ {\prime \prime} \longrightarrow \dots \longrightarrow \mathrm{R} _ {1} ^ {\prime \prime} \longrightarrow \mathrm{M} \longrightarrow 0 \\ 0 \longrightarrow \mathrm{N} \xrightarrow {f _ {n + 1} ^ {\prime}} \mathrm{R} _ {n} ^ {\prime} \longrightarrow \dots \longrightarrow \mathrm{R} _ {1} ^ {\prime} \xrightarrow {f _ {1} ^ {\prime}} \mathrm{M} \longrightarrow 0. \end{array}
\]

我们来证明 \(a\) )。设 \(\alpha \in \operatorname{Ext}_{\mathbf{A}}^{n}(\mathbf{M}, \mathbf{N})\) ，并设 \(\mathbf{P}\) 为 \(\mathbf{M}\) 的一个投射分解。设 \(a: \mathrm{P}(n) \to \mathbf{N}\) 为表示 \(\alpha\) 的复形态射；其唯一的非零分量是一个 A-同态 \(u: \mathrm{P}_{n} \to \mathbf{N}\) ，它满足 \(u \circ d_{n+1} = 0\) ，因此有分解 \(u = \overline{u} \circ \delta_{n}\) ，其中 \(\delta_{n}: \mathrm{P}_{n} \to \dot{\mathbf{Z}}_{n-1}\) 是由 \(d_{n}\) 诱导的映射（我们设 \(Z_{n-1} = \operatorname{Im} d_{n}\) ），而 \(\overline{u}\) 是从 \(Z_{n-1}\) 到 \(\mathbf{N}\) 的一个 A-同态。根据第 117 页的注记 1，类 \(\theta \in \operatorname{Ext}_{\mathbf{A}}^{n}(\mathbf{M}, Z_{n-1})\) ——即正合列

\[
0 \rightarrow Z _ {n - 1} \rightarrow P _ {n - 1} \rightarrow \dots \rightarrow P _ {0} \rightarrow M \rightarrow 0
\]

的类——等于态射 \((-1)^{n(n + 1) / 2}\delta_n\) 的同伦类。因此我们有

\[
\alpha = (- 1) ^ {n (n + 1) / 2} \overline {{{u}}} \circ \theta ,
\]

根据第 120 页的推论 3，由此可将 α 表示为某个正合列的类。

由 X 第 120 页的推论 1 可知，(ii) \(\Rightarrow\) (i) 且 (iii) \(\Rightarrow\) (i)。假设 (i) 成立，并设 P 为 M 的一个投射分解。存在一个交换图

\[
\begin{array}{c} 0 \longrightarrow \mathrm{N} \xrightarrow {f _ {n + 1}} \mathrm{R} _ {n} \longrightarrow \mathrm{R} _ {n - 1} \longrightarrow \dots \longrightarrow \mathrm{R} _ {1} \longrightarrow \mathrm{M} \longrightarrow 0 \\ u _ {n} \uparrow \quad \uparrow u _ {n - 1} \quad \uparrow u _ {n - 2} \quad \uparrow \quad \uparrow_ {\mathrm{I} _ {\mathrm{M}}} \\ \mathrm{P} _ {n} \xrightarrow {d _ {n}} \mathrm{P} _ {n - 1} \longrightarrow \mathrm{P} _ {n - 2} \longrightarrow \dots \longrightarrow \mathrm{P} _ {0} \longrightarrow \mathrm{M} \longrightarrow 0 \\ u _ {n} ^ {\prime} \downarrow \quad \downarrow u _ {n - 1} ^ {\prime} \quad \downarrow u _ {n - 2} ^ {\prime} \quad \downarrow \quad \downarrow_ {\mathrm{I} _ {\mathrm{M}}} \\ 0 \longrightarrow \mathrm{N} \xrightarrow {f _ {n + 1}} \mathrm{R} _ {n} ^ {\prime} \longrightarrow \mathrm{R} _ {n - 1} ^ {\prime} \longrightarrow \dots \longrightarrow \mathrm{R} _ {1} ^ {\prime} \longrightarrow \mathrm{M} \longrightarrow 0. \end{array}
\]

从 \(\mathrm{P}(n)\) 到 \(\mathbf{N}\) 的、由 \(u_{n}\) 与 \(u_{n}^{\prime}\) 定义的态射是同伦的，因为它们都属于类 \((-1)^{n(n + 1) / 2}\theta\) ，因此 \(u_{n}^{\prime} - u_{n}\) 具有 \(w\circ d_n\) 的形式，其中 \(w:\mathrm{P}_{n - 1}\to \mathrm{N}\) 是一个 A-同态。将 \(u_{n - 1}^{\prime}\) 换为 \(u_{n - 1}^{\prime} - f_{n + 1}^{\prime}\circ w\) ，并将 \(u_{n}^{\prime}\) 换为 \(u_{n}\) ，就可归结为 \(u_{n} = u_{n}^{\prime}\) 的情形。这使得我们可以构造一个新的具有正合行的交换图：

\[
\begin{array}{c c c c c c c} 0 \longrightarrow \mathrm{N} \longrightarrow \mathrm{R} _ {n} \longrightarrow \mathrm{R} _ {n - 1} \longrightarrow \dots \longrightarrow \mathrm{R} _ {1} \longrightarrow \mathrm{M} \longrightarrow 0 \\ 1 _ {\mathrm{N}} \uparrow & \uparrow v & \uparrow u _ {n - 2} & \uparrow & \uparrow^ {1 _ {\mathrm{M}}} \\ 0 \longrightarrow \mathrm{N} \longrightarrow \mathrm{N} ^ {\prime} \longrightarrow \mathrm{P} _ {n - 2}. \longrightarrow \dots \longrightarrow \mathrm{P} _ {0} \longrightarrow \mathrm{M} \longrightarrow 0 \\ 1 _ {\mathrm{N}} \downarrow & \downarrow v ^ {\prime} & \downarrow u _ {n - 2} ^ {\prime} & \downarrow & \downarrow^ {1 _ {\mathrm{M}}} \\ 0 \longrightarrow \mathrm{N} \longrightarrow \mathrm{R} _ {n} ^ {\prime} \longrightarrow \mathrm{R} _ {n - 1} ^ {\prime} \longrightarrow \dots \longrightarrow \mathrm{R} _ {1} ^ {\prime} \longrightarrow \mathrm{M} \longrightarrow 0 \end{array}
\]

其中 \(\mathbf{N}'\) 是 \(\mathbf{P}_{n-1} \oplus \mathbf{N}\) 除以由诸对 \((d_n(x), -u_n(x))\) （ \(x \in \mathbf{P}_n\) ）所构成的子模而得的商模，而 \(v\) (相应地， \(v'\) )是由映射 \(u_{n-1} \oplus f_{n+1}\) （相应地， \(u_{n-1}' \oplus f_{n+1}'\) ）通过取商定义的。因此条件 (ii) 得到满足。

再假设条件 (i) 成立，并设 E 为 N 的一个内射分解。存在一个交换图

\[
\begin{array}{c c c c c c c} 0 \longrightarrow \mathrm{N} \longrightarrow \mathrm{R} _ {n} \longrightarrow & \dots \longrightarrow \mathrm{R} _ {1} \longrightarrow \mathrm{M} \longrightarrow & 0 \\ 1 _ {\mathrm{N}} \Big \downarrow & \Big \downarrow v _ {0} & v _ {n - 1} \Big \downarrow & \Big \downarrow v _ {n} \\ 0 \longrightarrow \mathrm{N} ^ {\prime} \longrightarrow \mathrm{E} ^ {0} \stackrel {{\delta^ {0}}} {{\longrightarrow}} & \dots \longrightarrow \mathrm{E} ^ {n - 1} \stackrel {{\delta^ {n - 1}}} {{\longrightarrow}} \mathrm{E} ^ {n} \\ 1 _ {\mathrm{N}} \Big \uparrow & \Big \uparrow v _ {0} ^ {\prime} & v _ {n - 1} ^ {\prime} \Big \uparrow & \Big \uparrow v _ {n} ^ {\prime} \\ 0 \longrightarrow \mathrm{N} \longrightarrow \mathrm{R} _ {n} ^ {\prime} \longrightarrow & \dots \longrightarrow \mathrm{R} _ {1} ^ {\prime} \longrightarrow & \mathrm{M} \longrightarrow & 0 \end{array}
\]

并且如上面所示，可以假设 \(v_{n}^{\prime}=v_{n}\) 。于是我们有一个具有正合行的交换图

\[
\begin{array}{c c c c c c c c} 0 \longrightarrow N \longrightarrow R _ {n} \longrightarrow \dots \longrightarrow R _ {2} \longrightarrow R _ {1} \longrightarrow M \longrightarrow 0 \\ \downarrow^ {1 _ {N}} & \downarrow & & \downarrow & \downarrow & \downarrow^ {1 _ {M}} \\ 0 \longrightarrow N \longrightarrow E ^ {0} \longrightarrow \dots \longrightarrow E ^ {n - 2} \longrightarrow M ^ {\prime} \longrightarrow M \longrightarrow 0 \\ \uparrow^ {1 _ {N}} & \uparrow & & \uparrow & \uparrow & \uparrow^ {1 _ {M}} \\ 0 \longrightarrow N \longrightarrow R _ {n} ^ {\prime} \longrightarrow \dots \longrightarrow R _ {2} ^ {\prime} \longrightarrow R _ {1} ^ {\prime} \longrightarrow M \longrightarrow 0 \end{array}
\]

其中 \(\mathbf{M}' = \mathbf{M} \times_{\mathrm{R}_{1}} \mathbf{E}_{n-1}\) （参见 X，第 120 页，推论 3，b））。因此条件 (iii) 成立，这就完成了定理的证明。

注记 1. --- 若环 A 是诺特环，且 A-模 M 与 N 都是有限生成的，则由 a) 的证明可知，Ext \(_{A}^{n}\) (M, N) 的每个元素都是与某个正合列 0 → N → R \(_{n}\) → \ldots{} → R \(_{1}\) → M → 0 相伴的类，其中诸 R \(_{i}\) 都是有限生成的。

推论. --- 设 \(0 \to \mathbf{N} \xrightarrow{f} \mathbf{R} \xrightarrow{g} \mathbf{M} \to 0\) 与 \(0 \to \mathbf{N} \xrightarrow{f'} \mathbf{R}' \xrightarrow{g'} \mathbf{M} \to 0\) 为两个正合列， \(\theta\) 与 \(\theta'\) 为它们在 \(\operatorname{Ext}^1(\mathbf{M}, \mathbf{N})\) 中的相伴类。为使 \(\theta = \theta'\) ，必须且只需存在一个 A-同态 \(h: \mathbf{R} \to \mathbf{R}'\) 使得图表

\[
\begin{array}{c} \mathsf {R} \\ \mathsf {N} \\ \mathsf {f ^ {\prime}} \end{array} \begin{array}{c} \mathsf {f} \\ \mathsf {h} \\ \mathsf {R ^ {\prime}} \end{array} \begin{array}{c} \mathsf {g} \\ \mathsf {M} \\ \mathsf {g ^ {\prime}} \end{array}
\]

交换的。这样的同态必然是一个同构。

由命题 4 的推论 1，该条件是充分的。若 \(\theta = \theta'\) ，则我们有一个具有正合行的交换图：

\[
\begin{array}{c} 0 \longrightarrow \mathrm{N} \longrightarrow \mathrm{R} \longrightarrow \mathrm{M} \longrightarrow 0 \\ \uparrow_ {1 _ {\mathrm{N}}} \Big | _ {h ^ {\prime}} \Big | _ {1 _ {\mathrm{M}}} \\ 0 \longrightarrow \mathrm{N} \longrightarrow \mathrm{R} ^ {\prime \prime} \longrightarrow \mathrm{M} \longrightarrow 0 \\ \uparrow_ {1 _ {\mathrm{N}}} \Big | _ {h ^ {\prime \prime}} \Big | _ {1 _ {\mathrm{M}}} \\ 0 \longrightarrow \mathrm{N} \longrightarrow \mathrm{R} ^ {\prime} \longrightarrow \mathrm{M} \longrightarrow 0. \end{array}
\]

态射 \(h'\) 与 \(h''\) 由 X 第 7 页推论 3 知是同构，且 \(h = h'' \circ h'^{-1}\) 回答了该问题。最后一个断言由同上引文得出。

注记 2. --- 定理 1 把 Ext \(_{A}^{n}\) (M, N) 描述为正合列的等价类集合；通过结构转移，容易描述在此集合上得到的群法则。事实上，设 θ（相应地，θ′）为正合列 0 → N \(\xrightarrow{f_{n+1}}\) R \(_{n}\) \(\xrightarrow{f_n}\) \ldots{} → R \(_{1}\) → M → 0（相应地，0 → N \(\xrightarrow{f_{n+1}'}\) R \(_{n}'\) \(\xrightarrow{f_n'}\) \ldots{} → R \(_{1}'\) → M → 0）的类。设 Δ : M → M ⊕ M 与 ∇ : N ⊕ N → N 为 A-线性映射，由 Δ(x) = (x, x)（x ∈ M）与 ∇(y, z) = y + z（y, z ∈ N）定义。考虑映射

\[
m: \operatorname{Ext} _ {\mathbf {A}} (\mathbf {M}, \mathbf {N}) \oplus \operatorname{Ext} _ {\mathbf {A}} (\mathbf {M}, \mathbf {N}) \rightarrow \operatorname{Ext} _ {\mathbf {A}} (\mathbf {M} \oplus \mathbf {M}, \mathbf {N} \oplus \mathbf {N})
\]

是在第 121 页的注记中定义的。采用同上引文的记号，我们有 \(\nabla \circ i_{\mathbf{N}} = 1_{\mathbf{N}}\) 和 \(q_{\mathbf{M}} \circ \Delta = 1_{\mathbf{M}}\) ，从而 \(\theta + \theta' = \nabla \circ m(\theta, \theta') \circ \Delta\) 。鉴于同上引文及第 120 页的推论 3，这就给出一个类为 \(\theta + \theta'\) 的正合列：例如，若 \(n \geqslant 2\) ，可取序列

\[
0 \rightarrow \mathrm{N} \rightarrow \mathrm{R} _ {n} ^ {\prime \prime} \rightarrow \mathrm{R} _ {n - 1} \oplus \mathrm{R} _ {n - 1} ^ {\prime} \xrightarrow {f _ {n - 1} \oplus f _ {n - 1} ^ {\prime}} \dots \rightarrow \mathrm{R} _ {2} \oplus \mathrm{R} _ {2} ^ {\prime} \rightarrow \mathrm{R} _ {1} ^ {\prime \prime} \rightarrow \mathrm{M} \rightarrow 0
\]

其中 \(R_{n}^{\prime\prime}\) 是 \(R_{n} \oplus R_{n}^{\prime}\) 除以由诸对

\[
\left(f _ {n + 1} (x), - f _ {n + 1} ^ {\prime} (x)\right) \quad \text { for } x \in \mathbf {N},
\]

所构成的子模而得的商模，且其中 \(R_{1}^{\prime\prime}=R_{1}\times_{M}R_{1}^{\prime}\)
% semantic-fragments: legacy-input-seam
\relax{}
